\documentclass[11pt]{article}
\usepackage[margin=1.15in]{geometry}
\usepackage{amsmath,amssymb,amsthm,booktabs}
\usepackage[colorlinks,linkcolor=blue,citecolor=blue]{hyperref}

\theoremstyle{plain}
\newtheorem{theorem}{Theorem}
\newtheorem{lemma}[theorem]{Lemma}
\newtheorem{proposition}[theorem]{Proposition}
\theoremstyle{definition}
\newtheorem{definition}[theorem]{Definition}
\newtheorem{question}[theorem]{Question}
\newtheorem{remark}[theorem]{Remark}

\newcommand{\FF}{\mathbb F}
\newcommand{\ZZ}{\mathbb Z}
\newcommand{\QQ}{\mathbb Q}
\newcommand{\Cay}{\operatorname{Cay}}
\newcommand{\Spec}{\operatorname{Spec}}

\title{Two mechanisms for cospectral Cayley graphs\\ on the Heisenberg group}
\author{W.~Dahn}
\date{July 2026}

\begin{document}
\maketitle

\begin{abstract}
\noindent Let $H$ be the Heisenberg group of order $q^3$, $q$ an odd prime. An
\emph{inverse-closed section} assigns to each nonzero $v\in\FF_q^2$ a central
coordinate $c(v)$ with $c(-v)=-c(v)$, and determines a connected
$(q^2-1)$-regular Cayley graph $\Cay(H,S(c))$ on $q^3$ vertices. We report a
census at $q=5$: among $6000$ random sections, $66$ pairs are cospectral
without being related by the evident affine symmetries. Every one of them
arises by one of exactly two mechanisms, which we name and separate by an
integral invariant. In the first, the two Galois-conjugate ``sheets'' of the
spectrum are \emph{exchanged} between the square and nonsquare central
characters; in the second they \emph{coincide} sheet by sheet, so that the
pair carries identical spectral data at every central character and is
nonetheless inequivalent. The two mechanisms occur at strikingly similar
rates ($34$ and $32$), which we do not explain.
\end{abstract}

\noindent\textbf{2020 MSC.} 05C50, 05C25, 20C15, 05E30.\\
\textbf{Keywords.} Cospectral graphs, Cayley graphs, Heisenberg group,
critical group, Weyl representation.

\section{Setup}

Let $q$ be an odd prime and
\[
H=\{(a,b,c):a,b,c\in\FF_q\},\qquad
(a,b,c)(a',b',c')=(a+a',b+b',c+c'-ab'+a'b),
\]
with centre $Z=\{(0,0,c)\}$. A function
$c:\FF_q^2\setminus\{0\}\to\FF_q$ with $c(-v)=-c(v)$ is an \emph{inverse-closed
section}; it makes
$S(c)=\{(v,c(v)):v\neq0\}$ inverse-closed, so
$\Gamma(c):=\Cay\bigl(H,S(c)\bigr)$
is an undirected $(q^2-1)$-regular graph on $q^3$ vertices.

Two obvious symmetries act on sections. The group $\mathrm{GL}_2(\FF_q)$ acts by
$c\mapsto\det(g)\,c\circ g^{-1}$, and translations act by
$c\mapsto c+\ell$ for a linear form $\ell$; together they generate an
\emph{affine} action of order $|\mathrm{GL}_2(\FF_q)|\,q^2$ under which
$\Gamma(c)$ is unchanged up to isomorphism. Call two sections
\emph{equivalent} if they lie in one orbit. At $q=5$ the number of orbits is
$20{,}592$, so the action is very nearly free.

\begin{definition}
A pair of sections is \emph{genuine} if the graphs $\Gamma(c),\Gamma(c')$ are
cospectral but the sections are inequivalent.
\end{definition}

\subsection*{Sheets}

The spectrum of $\Gamma(c)$ decomposes along the characters of $Z$. The trivial
character contributes the principal eigenvalue $q^2-1$; each nontrivial
character $\psi_t$ contributes the spectrum of the $q\times q$ Weyl block
\[
B_t(c)=\sum_{v\neq0}\rho_t\bigl(v,c(v)\bigr),
\qquad
\rho_t(a,b,c)e_x=\zeta^{\,t(c+2xb+ab)}e_{x+a},
\]
each with multiplicity $q$. Blocks at $t$ and $t'$ are Galois conjugate when
$t/t'$ is a square, so at $q=5$ there are exactly two distinct
\emph{sheets},
\[
\Sigma^{+}(c)=\Spec B_1(c)\quad\text{(square coset)},\qquad
\Sigma^{-}(c)=\Spec B_2(c)\quad\text{(nonsquare coset)} .
\]
Cospectrality of $\Gamma(c)$ and $\Gamma(c')$ says only that the
\emph{multiset union} $\Sigma^{+}\cup\Sigma^{-}$ agrees; it does not say the
sheets agree individually. That is precisely the room in which the two
mechanisms live.

\section{The two mechanisms}

\begin{definition}
A genuine pair $(c,c')$ is
\begin{itemize}
\item a \emph{sheet exchange} if
$\Sigma^{+}(c)=\Sigma^{-}(c')$ and $\Sigma^{-}(c)=\Sigma^{+}(c')$ while
$\Sigma^{+}(c)\neq\Sigma^{-}(c)$;
\item a \emph{sheet coincidence} if
$\Sigma^{+}(c)=\Sigma^{+}(c')$ and $\Sigma^{-}(c)=\Sigma^{-}(c')$.
\end{itemize}
\end{definition}

A sheet exchange is cospectral because the union is symmetric in the two
sheets: the pair swaps which Galois conjugate is attached to which square
class of central characters, leaving the rational product unchanged. A sheet
coincidence is a strictly stronger coincidence --- the two graphs agree at
\emph{every} central character individually --- and is therefore invisible to
any invariant computed sheet by sheet.

\begin{proposition}[census]\label{prop:census}
Among $6000$ uniformly random inverse-closed sections at $q=5$ there are $66$
genuine pairs. Of these, $34$ are sheet exchanges and $32$ are sheet
coincidences. No genuine pair is of any other type.
\end{proposition}

\begin{proposition}[integral separation]\label{prop:smith}
Both members of every genuine pair have the same critical group. The critical
groups nonetheless distinguish the mechanisms: for the sheet-exchange pairs
examined the $5$-primary part is $(\ZZ/25)^{5}\oplus(\ZZ/5)^{16}$, while for
the sheet-coincidence pair it is $(\ZZ/25)^{15}\oplus(\ZZ/5)^{6}$.
\end{proposition}

So Smith torsion, which fails to separate the two graphs \emph{within} a pair,
does separate the two \emph{kinds} of pair. Weisfeiler--Leman refinement
separates the graphs within a pair (rank $19$ versus $18$ at the first
$2$-dimensional refinement), so no pair here is a counterexample to anything;
the content is the dichotomy, not any single graph.

\begin{remark}
Neither mechanism is a Godsil--McKay switch in the naive sense: an exhaustive
search over the $2^{10}$ partial applications of the exact offset difference
of one exchange pair, and over the natural central-fibre, maximal-abelian-coset
and central-slice vertex families, produces no switching set realizing the
target.
\end{remark}

\section{What we cannot explain}

\begin{question}\label{q:rates}
Why are the two mechanisms nearly equinumerous? The census gives $34$ and
$32$, well inside the fluctuation of a fair split, over a sample large enough
that a systematic factor of two would have shown. Is there a bijection between
sheet exchanges and sheet coincidences?
\end{question}

A natural candidate is the determinant twist $c\mapsto\det(g)\,c\circ g^{-1}$
with $\det g$ a nonsquare, which is an affine map and hence carries a section
to an equivalent one, but which exchanges the two sheets. Composing it with
one member of an exchange pair turns that pair into a coincidence pair of
\emph{equivalent} sections --- so the twist explains why the two phenomena are
related, but not why inequivalent representatives of each occur equally often.

\begin{question}
Do both mechanisms persist for all odd $q$, and with what densities? At $q=3$
there are no genuine pairs at all (the action has only two orbits), and at
$q=7$ the orbit count $1{,}939{,}395{,}416{,}499{,}131$ makes a random census
hopeless: a sample of $80$ sections yields $80$ distinct spectra. Only $q=5$
is presently within reach.
\end{question}

\subsection*{Verification}

All spectra are computed to $10^{-6}$ and compared as rounded multisets;
inequivalence is decided by an exhaustive search over the full affine group of
order $12{,}000$; critical groups are computed by exact $p$-adic elimination on
the reduced Laplacian. Programs and certificates are archived
in \cite{repo}.

\begin{thebibliography}{9}

\bibitem{BCN}
A.~E.~Brouwer, A.~M.~Cohen, and A.~Neumaier,
\emph{Distance-Regular Graphs},
Ergeb.\ Math.\ Grenzgeb.\ (3) \textbf{18}, Springer, Berlin, 1989.

\bibitem{GM}
C.~D.~Godsil and B.~D.~McKay,
\emph{Constructing cospectral graphs},
Aequationes Math.\ \textbf{25} (1982), 257--268.

\bibitem{Ma}
S.~L.~Ma,
\emph{A survey of partial difference sets},
Des.\ Codes Cryptogr.\ \textbf{4} (1994), 221--261.

\bibitem{PDSS}
J.~Polhill, J.~A.~Davis, K.~W.~Smith, and E.~Swartz,
\emph{Genuinely nonabelian partial difference sets},
J.~Combin.\ Des.\ \textbf{32} (2024), 351--370.

\bibitem{repo}
W.~Dahn, \emph{W33-Theory} source repository; witnesses
\texttt{w33\_pass456}, \texttt{w33\_pass479}--\texttt{w33\_pass482}.

\end{thebibliography}

\end{document}
